\section{Methodology}
\label{sec:methodology}

% This section is critical for acceptance. Reviewers will reject
% without a clear search and selection protocol.

\subsection{Search Protocol}

We conducted a systematic search across five sources:

\begin{enumerate}[leftmargin=*, nosep]
    \item \textbf{Academic databases:} arXiv (cs.AI, cs.CL, cs.IR,
    cs.MA), Semantic Scholar, Google Scholar. Search terms:
    ``LLM memory,'' ``agent memory,'' ``long-term memory language
    model,'' ``episodic semantic memory agent,'' ``forgetting LLM,''
    ``memory-augmented generation.''
    \item \textbf{Conference proceedings:} NeurIPS 2023--2025, ICML
    2024--2025, ICLR 2025--2026, ACL/EMNLP 2024--2025, AAAI 2024--2025,
    ACM UIST 2023.
    \item \textbf{GitHub:} Repositories with $>$500 stars tagged
    ``memory,'' ``agent memory,'' ``LLM memory,'' ``long-term memory.''
    \item \textbf{Practitioner sources:} Medium, DEV Community, Substack,
    Hacker News, X/Twitter threads with $>$100 engagements on AI
    agent memory.
    \item \textbf{Snowball sampling:} References cited by the six
    existing survey papers~\cite{sumers2024coala, gao2023rag,
    zhang2024survey, hu2025memory, wu2025memory, du2025memory}.
\end{enumerate}

\paragraph{Screening counts.}
Following PRISMA reporting guidelines~\cite{page2021prisma}, the initial
search retrieved approximately 520 candidate records across all five
sources. After deduplication, 387 unique systems were screened by title
and abstract. Of these, 249 proceeded to full-text review; 111 were
excluded at this stage (58 out-of-scope, 29 insufficient documentation,
24 duplicates of already-included systems). The final corpus comprises
138 systems. Figure~\ref{fig:prisma} visualizes the screening flow.
Table~\ref{tab:excluded} lists deliberate exclusions for systems that
required full-text review before exclusion. The complete system list
with per-system coding labels is in the system matrix
(Table~\ref{tab:system-matrix}).

\begin{figure*}[!t]
\centering
\begin{tikzpicture}[
  node distance=0.65cm and 2.0cm,
  mainbox/.style={rectangle, draw=black, thick, text width=5.8cm,
                  minimum height=1.5cm, align=center, font=\small},
  exclbox/.style={rectangle, draw=black!65, thick, text width=4.8cm,
                  minimum height=1.2cm, align=center, font=\footnotesize,
                  fill=gray!8},
  arr/.style={-{Stealth[length=3mm]}, thick},
  darr/.style={-{Stealth[length=3mm]}, thick, dashed, draw=black!60},
]
\node[mainbox] (id) {
  \textbf{Identification}\\[3pt]
  Records from 5 sources\\
  (arXiv/Semantic Scholar, conf.\ proceedings,\\
   GitHub, practitioner sources, snowball)\\[2pt]
  $n \approx 520$
};
\node[mainbox, below=of id] (screen) {
  \textbf{Screening}\\[3pt]
  Unique records after deduplication\\[2pt]
  $n = 387$
};
\node[mainbox, below=of screen] (elig) {
  \textbf{Eligibility}\\[3pt]
  Full-text assessed for inclusion\\[2pt]
  $n = 249$
};
\node[mainbox, below=of elig] (incl) {
  \textbf{Included}\\[3pt]
  Systems in final synthesis\\[2pt]
  $n = 138$
};
\node[exclbox, right=of screen] (screxcl) {
  Excluded at title/abstract screening\\
  (not LLM agent memory)\\
  $n = 138$\textsuperscript{$\dagger$}
};
\node[exclbox, right=of elig] (eligexcl) {
  Excluded after full-text review\\
  $n = 111$\\[3pt]
  \textit{Out of scope:}\ \ 58\\
  \textit{Insufficient docs:}\ 29\\
  \textit{Duplicate system:}\ 24
};
\draw[arr] (id) -- (screen);
\draw[arr] (screen) -- (elig);
\draw[arr] (elig) -- (incl);
\draw[darr] (screen.east) -- (screxcl.west);
\draw[darr] (elig.east) -- (eligexcl.west);
\end{tikzpicture}
\caption{PRISMA 2020 screening flow~\cite{page2021prisma}. Five sources were
searched; records were deduplicated, screened by title and abstract, then
assessed by full text.
\textsuperscript{$\dagger$}The coincidence $n = 138$ for both records
excluded at abstract screening ($387 - 249 = 138$) and systems included
in synthesis ($249 - 111 = 138$) is arithmetically exact but incidental.}
\label{fig:prisma}
\end{figure*}

\paragraph{Evidence-level markers.}
For commercial and closed-source systems where architectural details
cannot be independently verified from source code, claims in
\S\ref{sec:production} are annotated with evidence-quality superscripts:
\evD~official documentation, \evS~source code inspected,
\evV~vendor-reported (blog/press release),
\evI~inferred from behavior.
The full legend, assignment protocol, and reproducibility record are in
\Cref{sec:appendix-evidence}.

\subsection{Inclusion and Exclusion Criteria}

\begin{table}[htbp]
\centering
\caption{Inclusion and Exclusion Criteria}
\label{tab:criteria}
\begin{tabular}{@{}ll@{}}
\toprule
\textbf{Include} & \textbf{Exclude} \\
\midrule
Persistent memory across sessions & Session-only context management \\
Published 2023--2026 & Before 2023 (except seminal) \\
Usable system or reference impl. & Pure theoretical proposals \\
LLM-agent-focused & General database systems \\
\bottomrule
\end{tabular}
\end{table}

We note that multi-agent context management frameworks such as
Chain-of-Agents~\cite{zhang2024coa}, which partition long inputs across
collaborating workers within a single inference pass, fall on the
excluded side of this boundary: they address context-window limits
\emph{within} an interaction rather than maintaining persistent memory
\emph{across} interactions.

In addition to the criteria above, Table~\ref{tab:excluded} lists
specific systems that were evaluated during our search but deliberately
excluded because they fell outside our scope despite surface-level
relevance to agent memory.

\begin{table}[htbp]
\centering
\caption{Deliberately excluded systems with rationale.  Each system was
evaluated during our search but fell outside scope despite surface-level
relevance to agent memory.}
\label{tab:excluded}
\footnotesize
\setlength{\tabcolsep}{4pt}
\begin{tabular}{@{}p{0.33\columnwidth}p{0.60\columnwidth}@{}}
\toprule
\textbf{System} & \textbf{Reason for Exclusion} \\
\midrule
\multicolumn{2}{@{}l}{\textit{OpenClaw Ecosystem}} \\
TinyClaw        & $\sim$400~LOC bash wrapper; no novel memory \\
SecureClaw      & Security audit plugin; not a memory system \\
NanoClaw, ZeroClaw, NullClaw, PicoClaw & Language rewrites without novel memory mechanisms \\
\midrule
\multicolumn{2}{@{}l}{\textit{Agent Frameworks}} \\
BabyAGI         & Archived 2024; task-queue equivalent to AutoGPT (included) \\
MetaGPT         & Multi-agent orchestration; message passing only, no persistent memory \\
AgentVerse      & Multi-agent simulation; no cross-session memory \\
Autogen / AG2   & Multi-agent conversation framework; session-scoped message history only \\
SWE-agent       & Benchmark agent; no persistent memory across tasks \\
\midrule
\multicolumn{2}{@{}l}{\textit{Platforms \& Infrastructure}} \\
Cohere, Vectara & Search/retrieval platforms; infrastructure, not memory architectures \\
Semantic Kernel & SDK with memory connectors; infrastructure abstraction layer \\
\midrule
\multicolumn{2}{@{}l}{\textit{IDE / Coding Tools}} \\
Aider           & AI pair programmer; chat history only, no persistent memory architecture \\
Continue        & Open-source code assistant; context providers but no cross-session memory \\
Devin           & Autonomous engineer; memory architecture not published \\
\midrule
\multicolumn{2}{@{}l}{\textit{Libraries}} \\
MemEngine       & Reimplementation library; wraps existing backends with no novel design \\
\midrule
\multicolumn{2}{@{}l}{\textit{MCP Memory Servers (below threshold)}} \\
WhenMoon claude-memory-mcp & $\sim$57 stars; identity anchors without novel retrieval \\
Memory Anchor       & Too new; insufficient adoption data \\
Memorix             & Derivative; combines existing backends \\
\bottomrule
\end{tabular}
\end{table}

\subsection{Analysis Framework}

Each system is analyzed along four axes using three complementary
taxonomies:

\begin{itemize}[leftmargin=*, nosep]
    \item \textbf{CoALA}~\cite{sumers2024coala}: Memory modules
    (working, episodic, semantic, procedural), action space, decision
    procedure.
    \item \textbf{RAG Taxonomy}~\cite{gao2023rag}: Naive, Advanced,
    or Modular RAG.
    \item \textbf{Memory Operations}~\cite{zhang2024survey}: Writing,
    management, reading.
\end{itemize}


\subsection{Data Extraction and Coding}
\label{sec:coding}

For each system meeting our inclusion criteria, we performed structured
data extraction along four orthogonal axes. The author coded each system; a random 20\% sample ($n = 28$ systems) was independently
verified by a second coder (graduate research assistant), yielding
Cohen's $\kappa > 0.85$ on all axes. The four coding axes are:

\begin{enumerate}[leftmargin=*, nosep]
    \item \textbf{Structure.} Following CoALA~\cite{sumers2024coala}, we
    classified each system's memory modules as \emph{working} (in-context
    scratchpad), \emph{episodic} (timestamped interaction records),
    \emph{semantic} (extracted facts and generalizations),
    \emph{procedural} (reusable action sequences or skills), or
    \emph{implicit} (behavioral patterns encoded in retrieval weights or
    model parameters rather than explicit stores). A system may implement
    multiple types simultaneously.

    \item \textbf{Retrieval.} We mapped each system's retrieval pipeline
    to the RAG taxonomy of Gao~\etal~\cite{gao2023rag}. \emph{Naive RAG}
    denotes single-stage embed-and-retrieve over a flat index.
    \emph{Advanced RAG} adds pre-retrieval query rewriting, post-retrieval
    re-ranking, or both. \emph{Modular RAG} composes independently
    swappable retrieval components (\eg, sparse search, dense search,
    graph traversal, temporal filtering) via a routing or fusion layer.
    We additionally recorded the specific retrieval modalities employed
    (sparse/BM25, dense/vector, knowledge-graph traversal, recency decay,
    importance weighting) and fusion strategy where applicable.

    \item \textbf{Lifecycle.} Drawing on Zhang~\etal~\cite{zhang2024survey}
    and the memory operations taxonomy of Du~\etal~\cite{du2025memory}, we
    coded each system for four lifecycle operations: \emph{formation} (how
    new memories are created---verbatim storage, LLM extraction, or
    reinforcement-learned construction); \emph{consolidation} (whether raw
    memories are merged, deduplicated, or promoted across stores);
    \emph{forgetting} (whether the system implements any decay, eviction,
    or active deletion mechanism); and \emph{transformation} (whether
    memories change form over time, \eg, episodic records being
    generalized into semantic facts). Systems lacking a given operation
    were coded as ``none.''

    \item \textbf{Preference.} We assessed the mechanism, if any, by
    which user preferences and behavioral patterns emerge from accumulated
    memory. We distinguished four levels: \emph{none} (coded~N, no
    preference modeling); \emph{LLM-extracted} (coded~E, the system
    explicitly prompts an LLM to summarize or extract preferences from
    stored interactions); \emph{confidence-tracked} (coded~T, preferences
    are stored as belief objects with explicit confidence scores that
    evolve over interactions, requiring an LLM call at write time); and
    \emph{emergent} (coded~M, preferences arise implicitly through
    retrieval frequency, Hebbian-style strengthening, or reinforcement
    signals without any per-interaction extraction step). This axis
    extends beyond existing survey taxonomies and reflects the gap
    identified by Hu~\etal~\cite{hu2025memory} between static knowledge
    stores and adaptive personalization.
\end{enumerate}

This four-axis coding scheme is designed to be both comprehensive and
orthogonal: \emph{structure} describes \emph{what} is stored,
\emph{retrieval} describes \emph{how} it is accessed, \emph{lifecycle}
describes \emph{when and how} it evolves, and \emph{preference} describes
\emph{why} it matters for downstream agent behavior. The resulting
classifications are reported in the comparative matrices of
\Cref{sec:comparison}.

\subsection{Baseline Evaluation}
\label{sec:methodology-baseline}

In addition to the survey analysis, we conduct a controlled baseline
replication to provide an independent reference point for the published
benchmark scores compiled in \S\ref{sec:benchmark-cross-comparison}.  We
evaluate two canonical systems---full-context injection and na\"ive vector
RAG---on both LoCoMo (1{,}540 questions, categories~1--4) and LongMemEval\_S
(500 questions) using Gemini-2.0-Flash-001 as the generator and GPT-4o-mini as
the LLM judge, both accessed via OpenRouter.  Statistical significance is
assessed using McNemar's test~\cite{mcnemar1947} on paired questions, with
95\% Wilson score confidence intervals~\cite{wilson1927}.  Full experimental details, code, and raw results are
available at~\cite{alaya2026}.\footnote{\url{https://github.com/SecurityRonin/alaya/tree/main/bench}}

\subsection{Threats to Validity}
\label{sec:threats}

Several threats to validity constrain the generalizability of this survey.
We identify five principal concerns and the mitigations we applied.

\textbf{Recency bias.} The field of LLM agent memory is evolving at an
extraordinary pace. Our search covers publications and systems through
April~2026, but significant new architectures continue to appear on a
weekly basis. We mitigate this by anchoring our taxonomy in stable
cognitive-science frameworks (CoALA, complementary learning systems
theory) rather than system-specific features, so that the analytical
apparatus remains applicable to future systems even as the landscape
shifts.

\textbf{Selection bias.} Our GitHub search criteria (repositories with
$>$500 stars) systematically favor English-language, open-source projects
with strong community engagement. Closed-source production systems
(\eg, proprietary memory layers in commercial agent platforms) are
underrepresented because they publish limited technical details. We
partially address this through practitioner source monitoring (blog posts,
social media analyses, community tutorials) that surface implementation
details of systems not formally published in academic venues.

\textbf{Classification subjectivity.} Mapping heterogeneous systems onto
a unified taxonomy necessarily involves interpretive judgment. The CoALA
framework~\cite{sumers2024coala} was designed for cognitive architectures
with clear module boundaries, but many production systems blur these
distinctions---a single vector store may serve simultaneously as episodic
and semantic memory depending on how entries are queried. We mitigate this
through explicit operational definitions for each coding category
(Section~\ref{sec:coding}), independent dual coding, and transparent
reporting of borderline cases in our system analyses.

\textbf{Completeness.} No survey of this rapidly expanding field can
claim exhaustiveness. We identified over 100 systems meeting our inclusion
criteria, but new projects appear weekly on arXiv and GitHub. Rather than
pursuing completeness, we aimed for \emph{taxonomic saturation}: we
continued adding systems until new entries no longer introduced novel
architectural patterns or memory mechanisms. The final set is sufficient
to establish the taxonomy's axes and identify convergent design principles,
even if individual systems are inevitably omitted.

\textbf{Reproducibility.} Several systems in our corpus exist only as
blog posts, social media threads, or proprietary deployments with limited
technical documentation. For these systems, we relied on the available
descriptions and, where possible, examined open-source code. We flag
systems with limited documentation in our analyses and weight our
conclusions toward systems with reproducible implementations or
peer-reviewed publications. Where architectural details are ambiguous, we
report the ambiguity rather than speculate.
